\section{Related Work}
\label{sec:related}

\textbf{Model Merging and Task Vectors.} Model merging has emerged as a major focus for resource-efficient multi-task learning, allowing practitioners to combine the capabilities of multiple fine-tuned models without the overhead of joint training or multi-branch deployment. The seminal work of \citet{mitchell80} and subsequent developments in task arithmetic \cite{langley00, ilharco22} established that the difference between a fine-tuned model and its pre-trained base can be treated as a \emph{task vector} in parameter space. By linearly combining these task vectors, one can create a single model that exhibits multi-task functionality \cite{wortsman22}. Recent works have expanded this paradigm through more sophisticated static merging techniques, such as Fisher-weighted merging \cite{matena21}, RegMean (which aligns activations) \cite{jin22}, resolving interference \cite{yadav23}, permuting parameters \cite{ainsworth2022git}, and feature matching \cite{stoica2023zipit}. Despite their simplicity, static methods are fundamentally bounded because they enforce a static, uniform parameter configuration for all inputs. Consequently, they are highly prone to representation interference, where task-specific activations conflict and degrade overall performance.

\textbf{Dynamic Routing in Parameter Space.} To address representation interference, researchers have proposed dynamic model merging (also referred to as dynamic weight routing or on-the-fly model ensembling). Unlike standard token-level Mixture of Experts (MoE) \cite{shazeer2017outrageously, fedus2021switch} which routes tokens to different sub-networks, weight-space dynamic routing predicts input-dependent interpolation coefficients sample-by-sample to construct a custom merged model on-the-fly. This approach is highly efficient for low-rank adaptation (LoRA) settings \cite{hu2021lora}, where only a lightweight set of adapter weights is dynamically assembled, capping memory overhead at a negligible level. However, dynamic routing introduces a calibration phase: a small set of calibration samples is processed to optimize a lightweight parametric router. Our work addresses the fundamental generalization limits of this calibration process under severe data scarcity.

\textbf{Low-Data Calibration and Regularization.} Calibrating a dynamic router on a sparse dataset (e.g., $B_{\text{cal}} \le 64$) is a highly ill-posed problem. Standard parametric routing models with high capacity overfit the calibration split almost immediately, leading to a catastrophic drop in test-time accuracy—a phenomenon we refer to as generalization collapse. To combat this, several regularizers have been proposed. Task-Space Anchor Regularization (TSAR) anchors the router weights to predefined task centroids using proximal gradients. Task-Variance Regularization (VR-Router) forces the predicted routing coefficients to have low variance across a batch, preventing overconfident, single-task collapse. Additionally, non-parametric methods like Parameter-Free Subspace Routing (PFSR) bypass training altogether by utilizing raw representation similarities. While these heuristics offer practical improvements, they are fundamentally ad-hoc. They do not model the underlying geometry of the expert parameters, applying identical regularization penalties to all experts regardless of their parameter-space distances from the base model.

\textbf{Generalization Bounds and Rademacher Complexity.} Generalization theory seeks to bound the test-time performance of a model class using its training-time performance and a measure of the class's capacity. Rademacher complexity is a powerful and widely-adopted tool that measures a hypothesis class's ability to fit random noise \cite{bartlett2002rademacher, kearns89}. Classic results bound the Rademacher complexity of neural networks using their weight matrices' Frobenius or spectral norms \cite{kakade2008complexity, koltchinskii2002empirical, bousquet2002stability}, showing that generalization is closely tied to weight magnitudes. Furthermore, when composing neural network layers, Lipschitz continuity is often used to track the propagation of representation bounds \cite{mohri2018foundations, shalev2014understanding}. While these mathematical tools have been applied to standard supervised learning, they have not yet been used to analyze the generalization properties of dynamically merged models in weight space. Our work bridges this gap, deriving the first Rademacher-based generalization bound for weight-space dynamic routing and showing that parameter-space task-vector norms are the mathematically optimal scaling factors for routing weight decay.
